\documentclass[12pt]{article}
\usepackage[utf8]{inputenc}
\usepackage[T1]{fontenc}
\usepackage{hyperref}
\title{Gbe, Gbesi Language Resources for K-12 Curriculum Development}
\author{Compiled by Claude Sonnet 5.5 (claude-sonnet-5-5) for Andr\'as Kornai}
\date{}
\begin{document}
\maketitle

\section{Introduction}
Gbesi Gbe (ISO 639-3: gbs; Glottocode: gbes1238) is a Gbe language of southern Benin, in the Ayizo--Kotafon--Gbesi group. Joshua Project lists about 141,000--144,000 people. It is not known to be taught in schools. Very little was found: no grammar, dictionary or corpus appeared in the search results.

\section{Primary Educational Resources}
None were found.

\section{Linguistic and Reference Materials}
\begin{itemize}
\item PHOIBLE phonological inventory for Gbesi Gbe (47 segments): \url{https://phoible.org/languages/gbes1238}
\item Kluge (2011), a sociolinguistic survey of the Gbe communities of Benin and Togo (Glottolog record): \url{https://glottolog.org/resource/reference/id/55833}
\item Wikipedia article on the Ayizo language, which covers the group: \url{https://en.wikipedia.org/wiki/Ayizo_language}
\end{itemize}

\section{Religious and Community Resources}
\begin{itemize}
\item Joshua Project entry for Gbesi Gbe, listing Scripture status as translation started: \url{https://joshuaproject.net/languages/gbs}
\end{itemize}

\section{Digital and Audio Resources}
Joshua Project reports no audio Bible, Jesus Film or gospel recordings. No other audio or digital resources were found.

\section{Curriculum Development Context}
Almost nothing in the language was found. A curriculum would need primary documentation work first, drawing on related Gbe languages that have written traditions.
\end{document}